\subsection*{Part 1: Point in a direction (the table)}

\begin{center}
\renewcommand{\arraystretch}{1.3}
\begin{tabular}{P{2.8cm} P{6.0cm} P{3.0cm} P{2.8cm}}
\toprule
\textbf{Scratch direction} & \textbf{Python code} & \textbf{Python direction} & \textbf{Which way?} \\
\midrule
90 & \py|setheading(90 - 90)| & 0 & right \\
0 & \py|setheading(90 - 0)| & 90 & up \\
-90 & \py|setheading(90 - -90)| & 180 & left \\
45 & \py|setheading(90 - 45)| & 45 & up and right \\
135 & \py|setheading(90 - 135)| & $-45$ & down and right \\
\bottomrule
\end{tabular}
\end{center}

Taking away a minus number is the same as adding, so $90 - (-90) = 180$. Python understands $-45$ too: it is the same direction as 315.

\subsection*{Part 1: Translate it yourself}

\begin{code}
from ColabTurtlePlus.Turtle import *
clearscreen()
initializeTurtle((480, 360))
penup()
goto(-100, 0)
pendown()
pencolor("red")
pensize(4)
setheading(90 - 0)
forward(80)
right(90)
forward(80)
stamp()
\end{code}

The turtle draws a red line up from $(-100, 0)$, turns right, draws a line across, and stamps itself at $(-20, 80)$. Writing \py|setheading(90)| instead of \py|setheading(90 - 0)| also works. If you added \py|shape("turtle")|, well done: your stamp looks like a turtle.

\subsection*{Part 1: Two-line puzzle}

\begin{code}
setheading(90 - 180)
forward(100)
\end{code}

\py|setheading(-90)| works too, and so does \py|right(90)| instead of \py|setheading|, because the turtle starts facing right. A quarter turn clockwise makes it face down.
